\chapter{Conclusions} \label{cap:consideracoes}

In this monograph we presented the activities we developed in the course and some of the study materials we reviewed. With this project proposal for the Alware system, we studied topics in the intersection of safety analysis and embedded AI.

The literature review provided an encouraging outlook for the project's feasibility. Using TinyML architectures, high-accuracy driver state detection is achievable even on resource-constrained hardware \cite{essa2026} \cite{lumina2026}.

A significant portion of this study was dedicated to the project requirement specification and architecture. Moving forward into the Embedded Systems Laboratory course, our focus will shift from theoretical specification to practical development. This will include model training and optimization using TensorFlow Lite to compress deep learning models for the Raspberry Pi environment; physical integration of the camera to ensure good quality a real-time image data stream; and assembling the black box enclosure, ensuring the buzzer and LED alerts meet the specified intensity requirements for driver engagement.

In summary, while this is an initial study, the clarity of the requirements and the efficiency of the proposed architecture suggest that Alware is a feasible solution for enhancing driver safety through low-cost, high-efficiency embedded AI.